\cvsection{Expériences professionnelles}

\cvexperience{Été 2024 \\ (4 mois)}{Concepteur de circuits imprimés pour imagerie médicale à rayons X}
{Groupe de Recherche en Appareillage Médical de Sherbrooke}{Sherbrooke}
{Conception d'un \textit{PCB} 8 couches haute vitesse permettant la caractérisation de nouvelles méthodes de traitement de signaux analogiques provenant de détecteurs \textit{SiPM}}
{\vspace{-12pt}
\begin{itemize}
        \item Conception des premières simulations analogiques (\textit{SPICE}) du projet
        % \begin{itemize}
        %     \item Filtrage de signaux provenant d'expériences en laboratoire en Python
        %     \item Création, avec MATLAB, de modèles LTspice de pièces utilisées dans le circuit
        % \end{itemize}
        \item Élaboration du cahier des charges avec une équipe multidisciplinaire de chercheurs
        \item Utilisation des équipements d'assemblage et de tests des \textit{PCB}
        % \item Dessin de la schématique et du routage d'un PCB 8 couches haute vitesse
\end{itemize}}


\cvexperience{Aut. 2023 \\ Hiver 2023 \\ (8 mois)}{Développeur stack senseur quantique}{SBQuantum}{Sherbrooke}{Conception et développement de la chaîne d'acquisition et de traitement de données pour magnétomètres quantiques}
{\vspace{-12pt}
\begin{itemize}
    \item Implémentation du protocole \textit{UAV-CAN} pour la communication dans l'espace
    % \item Interfaçage d'un microcontrôleur ATSAME avec un FPGA iCE40
    % \item Migration de la chaîne d'outils Microchip vers des outils \textit{open source}
    % \begin{itemize}
    %     \item OpenOCD (programmation / communication JTAG)
    %     \item CMake / GNU Arm Embedded Toolchain (compilation)
    %     \item GDB / Cortex-Debug (débogage)
    % \end{itemize}
    \item Conception des \textit{PCB} formant la nouvelle génération de magnétomètres
    % \begin{itemize}
    %     \item Ajout d'une EEPROM permettant d'ajouter un segment de mémoire protégé
    %     \item LTSpice
    %     \item Altium Designer
    %     \item Écriture de pilotes en C
    % \end{itemize}
    % \item Programmation du \textit{Firmware Gold} et du \textit{Bootloader} pour assurer la sécurité du code une fois en orbite
    \item Programmation du \textit{Bootloader} pour assurer la sécurité du code une fois en orbite
    % \begin{itemize}
    %     \item Code assembleur AVR
    %     \item Programmation C \textit{bare-metal}
    %     \item GNU Linker Script Language
    % \end{itemize}
    % \item Écriture des pilotes en Python pour la station au sol
    \item Supervision et formation de stagiaires et de nouveaux employés
    % \item Rédaction de la documentation de formation
\end{itemize}}
    

\cvexperience{Été 2022 \\ (4 mois)}{Développeur de modules pour circuits reprogrammables (\textit{FPGA})}{Orthogone Technologies Inc.}{Montréal}{Développement de modules (\textit{IP cores}) de communication Ethernet à très basse latence}
{\vspace{-12pt}
\begin{itemize}
    % \item Interfaçage Ethernet vers AXI-Lite entre un CPU et un FPGA
    % \begin{itemize}
    %     \item Verilog
    %     \item Suite Xilinx (Vivado)
    % \end{itemize}
    % \item Mise en place des contraintes et des spécifications et écriture des bancs de test
    % \begin{itemize}
    %     \item SystemVerilog
    %     \item ModelSim et QuestaSim (Analyseur logique)
    % \end{itemize}
    % \item Amélioration du flux de travail avec Jenkins et Nexus Repository pour accélérer la conception et améliorer la sécurité des produits
    % \begin{itemize}
    %     \item Développement sur un serveur Linux CentOS 7
    %     \item Écriture de scripts Jenkins permettant la génération parallèle de \textit{bitstreams}
    %     \item Écriture de scripts Python générant des rapports de performances
    % \item Développement du système de stockage sécurisé des modules distribués aux clients avec Nexus Repository
    \item Développement d'un système de stockage sécurisé des modules distribués aux clients
    % \end{itemize}
\end{itemize}}


\cvexperience{Hiver 2021 \\ (4 mois)}{Développeur web \textit{full-stack}}{Lime Health Inc.}{Québec}{Conception d'une base de données pour une PME émergente visant à mieux gérer les données de patients québécois}
{\vspace{-12pt}
\begin{itemize}
    % \item Conception sécuritaire de l'architecture de la base de données
    % \begin{itemize}
    %     \item Choix des outils (Ruby on Rails)
    %     \item Sécurisation avec Auth0
    % \end{itemize}
    % \item Migration de l'interface de Vue.js vers React
    \item Démonstration du produit devant des médecins et des investisseurs du CHUS
\end{itemize}}
